\documentclass[12pt]{article}
\usepackage[utf8]{inputenc}
\usepackage[T1]{fontenc}
\usepackage{hyperref}
\title{Toba Language Resources for K-12 Curriculum Development}
\author{Compiled by Claude Sonnet 5 (claude-sonnet-5) for Andr\'as Kornai}
\date{}
\begin{document}
\maketitle

\section{Introduction}
Toba, also called Qom or Toba-Qom (ISO 639-3 tob), is a Guaicuruan language of the Gran Chaco in Argentina and Paraguay. Wikipedia cites 31,580 native speakers (2011), and other sources give about 19,810 in Argentina (2000). In 2010 the province of Chaco declared Qom one of four official provincial languages, alongside Spanish, Moqoit and Wich\'i. The sources report that Spanish is displacing Qom in daily communication.

\section{Primary Educational Resources}
In Chaco, bilingual intercultural schools use two teachers per classroom, one non-indigenous and one indigenous. Indigenous teachers now focus largely on revitalisation and on teaching Qom as a second language.
\begin{itemize}
\item Study of learning to read and write in bilingual Qom contexts (Rosemberg, Ojea and Alam, 2012): \url{https://l1research.org/article/download/27/25/49}
\item Ethnographic chapter on bilingual intercultural education in Chaco: \url{https://link.springer.com/chapter/10.1007/978-3-031-25584-7_15}
\item Article from an Argentine education journal (the search results gave no title): \url{https://fh.mdp.edu.ar/revistas/index.php/r_educ/article/download/9196/9489}
\end{itemize}

\section{Linguistic and Reference Materials}
\begin{itemize}
\item Wikipedia entry for Toba Qom, listing the dialects no'olxaxanaq, lañaxashec and tacshec: \url{https://wikipedia.com/wiki/ISO_639:tob}
\item Avellana (2012), Forma y Funci\'on, on Spanish--Toba contact: \url{https://revistas.unal.edu.co/index.php/formayfuncion/article/view/38504}
\item Revista de Dialectolog\'ia y Tradiciones Populares (June 2016) article: \url{https://doaj.org/article/ea7ace08db6c4de1a5e2c373d8cd7c0a}
\end{itemize}
The search reports the 1980 \emph{Vocabulario Toba-Qom} by Albert and Lois Buckwalter, and a new encyclopedic dictionary under way. I did not find a URL for either.

\section{Religious and Community Resources}
A Mennonite mission report describes a Toba-Qom team, working with the Argentine Bible Society, that completed a re-translation of the Old Testament and has started an encyclopedic dictionary: \url{https://www.mennonitemission.net/news/in-argentina-mission-includes-the-preservation-of-a-language/}

\section{Digital and Audio Resources}
The search mentions Hecht (2025) on the ``Voces Tobas'' YouTube channel, but gave no URL, so the channel itself is not listed.

\section{Curriculum Development Context}
Qom has official provincial status in Chaco and an existing bilingual intercultural education system, which is the natural base for curriculum work. Language shift among the young is the main concern. The search results give no evidence of a full native-language textbook series.
\end{document}